
\textbf{Note: This study uses proof-of-concept synthetic data to validate statistical methodology. All numerical results require empirical confirmation with real HuggingFace and GitHub data before publication.}

Despite widespread awareness of standardized ML dataset documentation frameworks (Gebru et al.'s Datasheets for Datasets: 3,142 citations; Mitchell et al.'s Model Cards: 2,899 citations), this study finds that only 7\% of sampled repositories achieve basic documentation completeness (DCS\_3 $\geq$ 2.4) within 90 days of initial release (95\% CI: [3.4\%, 13.8\%], N=100 synthetic samples, simulating HuggingFace repositories from 2022-2024 with $\geq 10$ stars). Using retrospective temporal measurement at T0+90 days via a 3-tier detection protocol (release tags $\rightarrow$ dataset commits $\rightarrow$ repository creation), this study establishes that the documentation gap exists from initial release rather than through post-release degradation. Repository commit velocity exhibits a strong positive correlation with documentation quality (Spearman $\rho$ = 0.951, p = 5.$32\times 10^{-52})$, while contributor count ($\rho$ = 0.028, p = 0.389) and issue responsiveness ($\rho$ = 0.061, p = 0.272) show no significant relationship. Licensing clarity is the weakest documentation component (27\% achieving $\geq 0.5$ score vs. 77\% for data context; $\chi^2$ = 24.04, p = 6.$03\times 10^{-6})$, with 73\% of sampled repositories lacking a LICENSE file. These findings suggest that documentation compliance may be driven by workflow integration during active development rather than by framework awareness or community size. The study proposes interventions targeting commit-triggered prompts and automated licensing templates.
